% concatenated from SL2RPaper.tex, intro.tex, background.tex, mainthm.tex, examples.tex, misc.tex
\documentclass{article}
\usepackage[utf8]{inputenc}
\usepackage[english]{babel}
\usepackage{blindtext}
\usepackage{amssymb}
\usepackage{amsthm}
\usepackage{amsmath}
\usepackage{tikz-cd}
\usepackage{setspace}
\usepackage{thmtools}
\usepackage{thm-restate}
\usepackage{hyperref}
\usepackage{cleveref}
\usepackage[margin=1.2in]{geometry}

\usepackage{graphicx}
\graphicspath{ {images/} }

\newcommand{\im}[1]{\mathrm{im} #1}
\newcommand{\gal}[1]{\mathrm{Gal} #1}
\newcommand{\aut}[1]{\mathrm{Aut} #1}
\newcommand{\inter}[1]{\mathrm{int} #1}

\theoremstyle{plain}
\newtheorem{theorem}{Theorem}[section]

\theoremstyle{definition}
\newtheorem{definition}[theorem]{Definition}

\theoremstyle{plain}
\newtheorem{prop}[theorem]{Proposition}

\theoremstyle{plain}
\newtheorem{corollary}[theorem]{Corollary}

\theoremstyle{plain}
\newtheorem{lemma}[theorem]{Lemma}

\theoremstyle{remark}
\newtheorem{remark}[theorem]{Remark}

\theoremstyle{remark}
\newtheorem{example}[theorem]{Example}

\theoremstyle{plain}
\newtheorem{conjecture}[theorem]{Conjecture}

\theoremstyle{plain}
\newtheorem{question}[theorem]{Question}

\declaretheorem[name=Theorem,numberwithin=section]{thm}

\setlength{\parindent}{0cm}
\onehalfspacing

\title{Detected Seifert surfaces and intervals of left-orderable surgeries}
\author{Yi Wang}
\date{}

\begin{document}
	
\maketitle
	
\begin{abstract}
Motivated by the $L$-space conjecture, we prove left-orderability of certain Dehn fillings on integral homology solid tori with techniques first appearing in the work of Culler-Dunfield \cite{cullerdunfield}. First, we use the author's previous results \cite{wang} to construct arcs of representations originating at ideal points detecting Seifert surfaces inside certain 3-manifolds. This, combined with the holonomy extension locus techniques of Gao \cite{gao}, proves that Dehn fillings near 0 of such 3-manifolds are left-orderable. We then explicitly verify the hypotheses of the main theorem for an infinite collection of odd pretzel knots, establishing previously unknown intervals of orderable Dehn fillings. This verifies the $L$-space conjecture for a new infinite family of closed 3-manifolds. 
\end{abstract}

\section{Introduction}

Given a 3-manifold $M$, it is an important problem to determine whether or not $\pi_1(M)$ can be ordered; for an exposition, see Clay-Rolfsen \cite{clayrolfsen}. This study is particularly motivated by the $L$-space conjecture of Boyer-Gordon-Watson \cite{bgw}. For an experimental treatment of this conjecture, see \cite{dunfieldcensus}. We have the following results on orderability of Dehn fillings of various knots, all of which are consistent with the $L$-space conjecture:
\begin{itemize}
	\item Zung \cite{zung} showed that all surgeries of the figure-eight knot have left-orderable fundamental group.
	\item Gao \cite{gaotwobridge}, Tran \cite{tran}, and Le \cite{khanh} have also established left-orderability for certain Dehn fillings on various two-bridge knots.
	\item Nie \cite{niepretzel} and Khan-Tran \cite{khantran} studied orderability of Dehn fillings of 3-strand pretzel knots.
	\item Hu \cite{hu}, Lin-Nie \cite{linnie}, and Turner \cite{turner} studied orderability of cyclic branched covers of various knot complements.
\end{itemize}   
Many of the techniques to prove left-orderability involve representations into $\widetilde{PSL_2(\mathbb{R})}$, which is the only orderable Lie group. This approach has been systematically developed by Culler-Dunfield \cite{cullerdunfield} and Gao \cite{gao}. This paper utilizes the approach of \cite{gao}, which involves the \emph{holonomy extension locus}. 

\medskip

The general approach of \cite{cullerdunfield} and \cite{gao} is to identify smooth points on the $SL_2(\mathbb{C})$ character variety which may be conjugated to real representations, and utilize the smoothness property to deform these points into arcs of $SL_2(\mathbb{R})$-representations, which then lift to arcs of $\widetilde{PSL_2(\mathbb{R})}$-representations. These arcs may then correspond to intervals of orderable Dehn fillings. Such ``base" representations have included real Galois conjugates of discrete faithful representations into $SL_2(\mathbb{C})$, hyperbolic Dehn filling points, or points coming from roots of the Alexander polynomial \cite{cullerdunfield} \cite{gao}.

\medskip

The central idea of this paper is to deform $SL_2(\mathbb{R})$-representations from \emph{ideal points} on the character variety whose \emph{limiting characters} are known to be characters of $SL_2(\mathbb{R})$-representations; this knowledge comes from the topology of surfaces associated to these ideal points via Culler-Shalen theory \cite{cullershalen}. Our results utilize an understanding of the limiting character developed in the work of Paoluzzi-Porti \cite{paoluzziporti}, Tillmann \cite{tillmann}, and the author \cite{wang}. Ideal points are a new potential source of $SL_2(\mathbb{R})$ represenations of one-cusped 3-manifolds that can exploited to prove left-orderability of their Dehn fillings through the holonomy extension locus techniques of \cite{gao}.

\begin{definition}
We say that $M$ is \emph{half-orderable near 0} if there exists some $a \in \mathbb{R}$ such that $M(r)$ for $r \in (0, a)$ has a left-orderable fundamental group. (If $a < 0$, replace $(0, a)$ with $(a, 0)$.)
\end{definition}

The following theorem follows quickly from the techniques of \cite{gao}, and motivates the construction of $SL_2(\mathbb{R})$-representations coming from ideal points on character varieties. See Section \ref{sec:background} for a complete explanation of the terminology used in the theorem statement.

\begin{theorem}
Let $M$ be an integral homology solid torus with a genus $g$ Seifert surface $S_g$. Suppose that $S_g$ is detected by an ideal point $x$ on a component of the character variety on which the trace of the longitude is non-constant, and there is an arc of representations $\rho_t: \pi_1(M) \to SL_2(\mathbb{R})$ limiting toward $x$ such that for any lift $\widetilde{\rho_t}: \pi_1(M) \to \widetilde{PSL_2(\mathbb{R})}$, the translation number of $\widetilde{\rho_t}(\ell)$ is 0. Then $M$ is half-orderable near 0. 
\end{theorem}

This theorem offers a systematic approach to explaining particular arcs appearing in the holonomy extension locus of \cite{gao}: those with horizontal asymptotes at the $x$-axis. The setting of Seifert surfaces and boundary slope 0 fits well within our framework for several reasons. First, the author's previous results \cite{wang} already provide many suitable candidates for which the ideal point hypothesis is satisfied. Second, for ideal points detecting Seifert surfaces, the limiting character is well-understood enough so that the translation number hypothesis is known to be satisfied often. Thus, in this setting, the only obstruction to proving half-orderability near 0 is finding arcs of $SL_2(\mathbb{R})$-representations limiting toward ideal points. The construction of such arcs forms the technical core of the paper. For an explicit infinite family of knots for which all the hypotheses are satisfied, we prove the following corollary:

\begin{corollary}\label{thm:pretzel}
The $(-3, 3, 2n+1)$ pretzel knot complements are half-orderable near 0.
\end{corollary}

This is a new explicit addition to list of left-orderable 3-manifolds. The $(2p+1, 2q+1, 2r+1)$ pretzel knots with $p, q, r > 0$ were already known to be orderable in a neighborhood of 0 from the work of Khan-Tran \cite{khantran}, which proved using elliptic representations that such knots are orderable in the interval $(-\infty, 1)$. In contrast, the results in this paper utilize arcs of boundary-hyperbolic representations into $SL_2(\mathbb{R})$ ending at ideal points detecting Seifert surfaces, establishing the presence of certain arcs in the holonomy extension locus. The process through which Corollary \ref{thm:pretzel} is proven demonstrates a general methodology for proving similar results for other knot complements with detected Seifert surfaces. 

\medskip

Corollary \ref{thm:pretzel} also proves the $L$-space conjecture for Dehn fillings near 0 of these pretzel knots, after combining with the results of Li-Roberts \cite{liroberts} and Hedden \cite{hedden}. 

\subsection{Outline of the paper}

Section \ref{sec:background} goes through necessary background to understand the arguments of the paper, including character varieties, Culler-Shalen theory, left-orderability, and related techniques. Section \ref{sec:proof} contains the main results of the paper. Section \ref{sec:examples} explicitly verifies the hypotheses of the main theorem for a collection of odd pretzel knots. Section \ref{sec:future} lists some potential future directions of research that would improve the scope and applicability of Theorem \ref{thm:main}. 

\subsection{Acknowledgements}

The author thanks Nathan Dunfield for numerous helpful conversations that guided him through this work. 

\section{Background}\label{sec:background}

For relevant background on 3-manifolds, JSJ decompositions, and Seifert fiberings, see \cite{thurston}. For more background on Culler-Shalen theory and limiting characters, see the background section of \cite{wang}, and for background on left-orderability and the holonomy extension locus, see \cite{gao}.

\subsection{Ideal points and incompressible surfaces}

Let $M$ be a compact 3-manifold. 

\begin{definition}
The \emph{$SL_2(\mathbb{C})$ character variety} of $M$, denoted $X(M)$, is defined as 
\begin{equation}
	X(M) = \{\text{tr}(\rho) \mid \rho: \pi_1(M) \to SL_2(\mathbb{C})\}
\end{equation}
i.e. the set of traces of representations to $SL_2(\mathbb{C})$.
\end{definition}

One can see that the $SL_2(\mathbb{C})$ is an affine algebraic set, often with multiple algebraically irreducible components. See \cite{cullershalen} for a proof. 

\begin{remark}
Note that this is \emph{not} the same as the moduli space of conjugacy classes of $SL_2(\mathbb{C})$ representations. In particular, these moduli spaces can exhibit non-Hausdorff topology near conjugacy classes of reducible representations. 
\end{remark}

\begin{definition}
Let $M$ be a hyperbolic 3-manifold. The \emph{holonomy representation} $\rho_0: \pi_1(M) \to SL_2(\mathbb{C})$ is the unique discrete faithful representation corresponding to the hyperbolic structure of $M$. If $M$ is Seifert-fibred over a hyperbolic 2-orbifold, a \emph{holonomy representation} is equal to $\pm I$ on the regular fiber, and equal to the discrete faithful representation on the fundamental group of the base orbifold. A \emph{canonical component} is an algebraically irreducible component which contains the trace of a Galois conjugate of a holonomy representation. 
\end{definition}

A foundational result of Thurston \cite{thurston} is that the complex dimension of the canonical component is equal to the number of cusps in a hyperbolic 3-manifold. One particular focus of this paper will be \emph{ideal points}, which can informally be thought of as ``points at infinity" on affine algebraic curves. 

\begin{definition}
Let $C$ be an affine algebraic curve, and let $\widetilde{C}$ be a projectivization. An \emph{ideal point} of $C$ is an element $x \in \widetilde{C} \setminus C$. 
\end{definition}
%
%\begin{definition}
%Let $M$ be an integral homology solid torus. We say that $M$ is \emph{longitudinally rigid} if $M(0)$ is 0-dimensional.
%\end{definition}
%
%Longitudinal rigidity is an important hypothesis in the results of \cite{cullerdunfield} and \cite{gao}. For this paper, we will need a similar, but weaker condition. 
%
%\begin{definition}	
%Let $S$ be an incompressible surface in $M$ with boundary $\ell$, and let $r: X(M) \to X(M \setminus S)$ be the map induced by restriction. We say that a cusped hyperbolic 3-manifold is \emph{longitudinally semi-rigid} if the image of $r$ intersected with the variety setting the trace of $r(\ell)$ equal to 2 is 0-dimensional. 
%\end{definition}
%
%It is clear that longitudinal rigidity implies longitudinal semi-rigidity. We will see in Section \ref{sec:examples} that there exist one-cusped 3-manifolds which are longitudinally semi-rigid but not longitudinally rigid. The proof of Theorem \ref{thm:main} only needs longitudinal semi-rigidity.

%\subsection{Sutured gluing subvarieties}
%
%Let $M$ be a 3-manifold such that $\partial M$ is an incompressible genus $2g$ surface, denoted $S_{2g}$. Let $S_{2g} = P^+ \cup c \cup P^-$ be a suturing of the genus $2g$ surface, where $P^{\pm}$ is a once-punctured genus $g$ surface, glued among the separating simple closed curve $c$. (\textbf{ADD IMAGE}) Pick an orientation-preserving automorphism $\varphi: P^+ \to P^-$, and let $M(\varphi)$ be the 3-manifold obtained by gluing $P^+$ to $P^-$ via $\varphi$. Then $M(\varphi)$ is a manifold with one torus boundary component. Assume that $M(\varphi)$ is either hyperbolic or Seifert-fibered over a hyperbolic 2-orbifold, and let $\rho_{\varphi}$ be the holonomy representation of $M(\varphi)$. 
%
%\begin{definition}
%Let $\pi_1(P^\pm) \cong F_{2g} = \langle a_1, \dots, a_{2g} \rangle$. Let $w_1, \dots, w_{6g - 6}$ be the set of words $a_i$, $1 \leq i \leq 2g$ and $a_{i_1}a_{i_2}$, where $1 \leq i_1 < i_2 \leq 2g$. Define the \emph{sutured gluing subvariety} $X(M, P^+, P^-) \subset X(M)$ to be the subvariety cut out by the equations $\chi(w_i) = \chi(\varphi_*(w_i))$. 
%\end{definition}
%
%\begin{remark}
%The character of an $F_{2g}$ representation is determined entirely by $w_1, \dots, w_{6g - 6}$; see \textbf{CITE ASHLEY / GOLDMAN}.
%\end{remark}
%
%The sutured gluing subvariety captures exactly the $SL_2(\mathbb{C})$-representations of $M$ which extend to representations of $M(\varphi)$, whose fundamental group is an $HNN$-extension of $\pi_1(M)$. If $r: X(M(\varphi)) \to X(M)$ is the map induced by restriction, one can see that $X(M, P^+, P^-)$ contains $r(\chi(\rho_\varphi))$. Let $\rho^* \in X(M, P^+, P^-) \setminus r(X(M(\varphi)))$ be an isolated point. 
%
%\begin{definition}
%We say that $M$ is \emph{smoothly sutured at $\rho^*$} if the trace of $\rho^*$ is a smooth point in $X(M, P^+, P^-)$. 
%\end{definition}
%
%\begin{definition}
%Let $\pi_1(M) = \langle b_1, \dots, b_n \mid r_1, \dots, r_m \rangle$, with $\partial M = S_{2g} = P^+ \cup c \cup P^-$ a suturing of the boundary, and $\pi_1(P^+) = \langle a_1, \dots, a_{2g} \rangle$, where $a_i$ is a word in the $b_i$. Let $\varphi: P^+ \to P^-$ be an orientation-reversing homeomorphism. The \emph{Lin presentation} pf $\pi_1(M(\varphi))$ is 
%\begin{equation}
%	\pi_1(M(\varphi)) = \langle b_1, \dots, b_n, t \mid r_1, \dots, r_m, ta_it^{-1} = \varphi_*(a_i) \rangle
%\end{equation}
%\end{definition}
%
%This is essentially writing down $\pi_1(M(\varphi))$ as an $HNN$-extension of $\pi_1(M)$. Using the $SL_2(\mathbb{C})$ trace relation
%\begin{equation}
%	\text{tr}(XY) = \text{tr}(X)\text{tr}(Y) - \text{tr}(XY^{-1})
%\end{equation}
%we can express $\text{tr}(a_i)$ as polynomials in $\text{tr}(b_i)$ and $\text{tr}(b_ib_j), 1 \leq i < j \leq n$. The sutured gluing subvariety of $M(\varphi)$ is therefore cut out by the equations
%\begin{equation}
%	X(M, P^+, P^-) = \{\text{tr}(a_i) = \text{tr}(\varphi_*(a_i))\} \subset X(F_{2g})
%\end{equation} 
%We make the following observation.
%
%\begin{corollary}\label{cor:defined}
%The smooth suturing subvariety is defined over $\mathbb{Q}$. 
%\end{corollary}
%
%The Lin presentation will be most useful in the case where $M$ is a genus 2 handlebody, since $X(M) = \mathbb{C}^3$, and we will be dealing with explicit equations in $\mathbb{C}^3$. Indeed, if we write this down as $F_2 \cong \langle a, b \rangle$, we can write $x = \text{tr}(a), y = \text{tr}(b), z = \text{tr}(ab)$. This notation will be used throughout Section \ref{sec:examples}. 

In the original paper introducing $SL_2(\mathbb{C})$-character varieties in the context of 3-manifolds, Culler and Shalen established a method which associates ideal points on a curve in $X(M)$ to incompressible surfaces in 3-manifolds. The theory is developed in full detail in \cite{cullershalen} and Chapter 19 of \cite{handbook}, while brief outlines can be found in \cite{paoluzziporti} and \cite{wang}.

\begin{definition}
Let $M$ be a hyperbolic 3-manifold, and let $X(M)$ be the $SL_2(\mathbb{C})$ character variety. If an incompressible surface $S \subset M$ is associated to an ideal point $x$ of $X(M)$, we say that $S$ is \emph{detected} by $x$. 
\end{definition}

Determining which incompressible surfaces in a given 3-manifold $M$ are detected by ideal points in $X(M)$ is a subtle problem, which has been studied in \cite{Cooper1994PlaneCA}, \cite{paoluzziporti}, \cite{Chesebro2005NotAB}, for instance. Some developments in Culler-Shalen theory, initiated by Dunfield \cite{Dunfield2008AP1} and Tillmann \cite{tillmann} have centered on the \emph{limiting character} at $x$. See \cite{wang} for a detailed exposition of the limiting character and its relevance to Culler-Shalen theory.

\begin{definition}
Let $x$ be an ideal point in $X(M)$ with detected surface $S \subset M$. Denote $M_i$ to be the connected components of $M \setminus S$. Let $\{\chi_j\} \subset X(M)$ be a sequence approaching $x$. The \emph{limiting character} $\chi_{\infty, i}: \pi_1(M) \to \mathbb{C}$ is given by 
\begin{equation}
	\chi_{\infty, i} = \lim_{\chi_j \to x}\chi_j|_{\pi_1(M_i)}
\end{equation}
It is shown in \cite{cullershalen} that $\chi_{\infty, i}$ takes on finite values on $\pi_1(M_i)$, and that $\chi_{\infty, i}$ is the trace of a well-defined representation $\rho_{\infty, i}: \pi_1(M_i) \to SL_2(\mathbb{C})$.
\end{definition}

In \cite{wang} and \cite{twicepunctured} the author initiated the study of limiting characters of punctured tori in hyperbolic knot complements. Limiting characters at ideal points detecting Conway spheres were also discussed in \cite{paoluzziporti}. We now state the relevant results on the limiting character that will be used for this paper. 

%We first define a technical consideration.
%
%\begin{definition}
%Let $M$ be a closed 3-manifold with one JSJ torus with JSJ component $M'$, and the gluing map between tori on $M'$ takes $\langle m_1, \ell_1 \rangle \mapsto \langle m_2, \ell_2 \rangle$. Let $\rho$ be a holonomy representation of $M'$, and let $\tau_i(\rho)$ be the ratio of the upper-right entries of $\rho(m_i)$ and $\rho(\ell_i)$ when conjugating $\langle m_i, \ell_i \rangle$ to an upper-triangular parabolic subgroup. We say $M$ is \emph{good} if the JSJ component is either Seifert-fibered or hyperbolic such that $\rho$ admits a Galois conjugate $\rho': \pi_1(M') \to SL_2(\mathbb{R})$ such that $\tau_1(\rho') \neq -\tau_2(\rho')$. If $M$ has multiple JSJ tori, we say that $M$ is \emph{good} if all gluings satisfy a similar property.
%\end{definition}
%
%This is to rule out a highly unlikely, degenerate case. 

The arguments of \cite{wang} can be adapted to prove the following, slightly modified theorem. 

\begin{theorem}[\cite{wang}]\label{thm:genusone}
Let $M$ be a one-cusped hyperbolic 3-manifold, and let $\mathcal{T} = \{T_i\}_{i=1}^n \subset M$ be a system of disjoint non-parallel, non-fibered genus one Seifert surfaces that cap off to a JSJ decomposition of $M(0)$. Let $r: X(M) \to X(M \setminus \mathcal{T})$ be the restriction map, and let $V \subset X(M \setminus \mathcal{T})$ be the Zariski-closure of the image of $r$. Suppose that the holonomy traces of the JSJ complementary regions of $M(0)$ are a point $\chi \in V$ such that the trace of the longitude is non-constant near $\chi$. Then $\mathcal{T}$ is detected by an ideal point $x$ on an irreducible component of $X(M)$ for which the trace of the longitude is nonconstant. In addition, the limiting character at $x$ is the trace of the holonomy representation of the JSJ complementary regions of $M(0)$. 
\end{theorem}

This is essentially a rephrasing of the theorem with extra hypotheses to allow for broader scope; the proof is exactly the same as the one presented in \cite{wang}. We now sketch the ideas in this proof. 
%We make the following remarks explaining the differences between this stated version of the theorem and the original version in \cite{wang}. 
%
%\begin{remark}
%The version of the theorem stated in \cite{wang} has different hypotheses. However, with the hypotheses stated in this version of the theorem, the exact same proof goes through. See \cite{twicepunctured} for a proof of an analogous theorem utilizing these hypotheses.
%\end{remark}
%
%\begin{remark}
%The original version of the theorem also excluded the case of hyperbolic manifolds on two sides of JSJ tori. This case can be proven with the addition of the following lemma.
%\end{remark}
%
%\begin{lemma}
%Suppose we have a circular arrangement of two-cusped hyperbolic 3-manifolds $M_i$ glued to each other to form a closed 3-manifold $M$. Let $\rho_i: \pi_1(M_i) \to SL_2(\mathbb{C})$ be a lift of the holonomy representation of $M_i$, chosen in such a way to be $SL_2(\mathbb{C})$-compatible. Then there does not exist a representation $\rho: M \to SL_2(\mathbb{C})$ which restricts to $\rho_i$ for all $i$.
%\end{lemma}
%
%\begin{proof}
%If $\widehat{T_k}$ bounds two hyperbolic spaces $M_1$ and $M_2$, let the fundamental group of the copies of the boundary torus be given by $\langle m_i, \ell_i \rangle \subset \pi_1(\partial M_i)$. Let $\rho_1$ and $\rho_2$ be the hyperbolic holonomy representations, so 
%\begin{equation}
%	\rho_i(m_i) = \pm\begin{pmatrix}1 & 1 \\ 0 & 1\end{pmatrix} \ \ \ \ \ \rho_i(\ell_i) = \pm\begin{pmatrix}1 & \tau_i \\ 0 & 1\end{pmatrix}
%\end{equation}
%We have that $\varphi$ sends $\ell_i$ to \textbf{USE ORIENTATION REVERSING, FINISH}.
%\end{proof}

For the purposes of this paper, we focus on the case where the Seifert surface is connected. 
\begin{itemize}
	\item Cut $M$ along the genus one Seifert surface $S$ to get a sutured 3-manifold $H = M \setminus S$ with genus 2 boundary $\partial H = P^+ \cup c \cup P^-$, where $P^\pm$ are once-punctured tori that are copies of $S$. Define the restriction map $r: X(M) \to X(H)$.
	\item Notice that $H$ with a 2-handle glued to $c$ is homeomorphic to the JSJ component of $M(0)$, equipped with a holonomy representation $\rho^*$.
	\item Let $C$ be the Zariski-closure of the image of $r$. Use the topology of the JSJ decomposition to determine that the character of $\rho^*$ is a point in $C$, and it is an isolated point in $C \setminus \text{im}(r)$. 
	\item Use Tillmann's result \cite{tillmann} characterizing ideal points in terms of limiting characters to conclude that $r^{-1}(\chi(\rho^*))$ is actually an ideal point of $X(M)$ which detects $S$, with limiting character $\chi(\rho*)$. 
\end{itemize}
The central idea of this paper is to take cases where the characters of $\rho^*$ have entirely real characters, deform to an analytic arc of real points in the image of $r$, and pull back to an arc of $SL_2(\mathbb{R})$-representations in $X(M)$. The above proof then relates this to an ideal point with real limiting character which detects a Seifert surface. This will have practical implications for proving certain Dehn fillings of $M$ are \emph{left-orderable}. 

\subsection{Left-orderability}

\begin{definition}
Let $G$ be a finitely generated group. We say $G$ is \emph{left-orderable} if it admits an ordering $<$ such that for all $f, g, h \in G$, $g < h$ implies $fg < fh$. If $M$ is a 3-manifold, we say that $M$ is \emph{left-orderable} if $\pi_1(M)$ is a left-orderable group.
\end{definition}

The interest in left-orderability in the context of low-dimensional topology comes from the \emph{$L$-space conjecture} from \cite{bgw}:

\begin{conjecture}
Let $M$ be a 3-manifold whose homology groups are isomorphic to those of $S^3$. Then the following are equivalent:
\begin{enumerate}
	\item $M$ is not an $L$-space.
	\item $M$ admits a taut foliation.
	\item $\pi_1(M)$ is left-orderable.
\end{enumerate}
\end{conjecture}

In general, left-orderability has been the most difficult of the three conditions to study. Here, we recall some of the existing literature on left-orderability of closed 3-manifold groups, with a particular emphasis on the techniques involving the translation extension locus and the holonomy extension locus. In \cite{cullerdunfield}, a systematic way to study left-orderability of 3-manifolds using character varieties was developed. It was shown in \cite{wiest} that if $H$ is a left-orderable group, any 3-manifold group $G$ with a homomorphism $G \to H$ is also orderable. The overarching strategy revolves around $\widetilde{PSL_2(\mathbb{R})}$, the universal cover of $SL_2(\mathbb{R})$ and $PSL_2(\mathbb{R})$. 

\begin{lemma}\cite{ghys}
$\widetilde{PSL_2(\mathbb{R})}$ is left-orderable.
\end{lemma}

The strategy developed by \cite{cullerdunfield} utilizes $SL_2(\mathbb{C})$ character varieties of one-cusped hyperbolic 3-manifolds to find homomorphisms from fundamental groups of Dehn fillings to $\widetilde{PSL_2(\mathbb{R})}$, showing that these Dehn fillings are also left-orderable. From $SL_2(\mathbb{C})$ character varieties, real points may correspond to $SL_2(\mathbb{R})$ representations. Given such a representation $\rho: \pi_1(M) \to SL_2(\mathbb{R})$, the \emph{Euler class} $e(\rho) \in H^2(\pi_1(M); \mathbb{Z})$ serves as an obstruction from lifting this representation to $\widetilde{PSL_2(\mathbb{R})}$. See \cite{ghys} for the definition, and \cite{cullerdunfield} for its role as an obstruction to lifting representations. For the purposes of this paper, we will restrict our attention to the case where $M$ is an integral homology solid torus, meaning that $H^2(\pi_1(M); \mathbb{Z})$ is trivial, and all $SL_2(\mathbb{R})$ representations will lift. 

\begin{definition}
The \emph{translation number} of a group element $g \in \widetilde{PSL_2(\mathbb{R})}$ is defined as 
\begin{equation}
	\text{trans}(g) = \lim_{n \to \infty}\frac{g^n(x) - x}{n} \ \ \ \ \ x \in \mathbb{R}
\end{equation}
This is independent of the choice of real number $x$.
\end{definition}

\begin{definition}
A matrix $A \in SL_2(\mathbb{R})$ is:
\begin{enumerate}
	\item \emph{elliptic} if $|\text{tr}(A)| < 2$
	\item \emph{parabolic} if $|\text{tr}(A)| = 2$
	\item \emph{hyperbolic} if $|\text{tr}(A)| > 2$
\end{enumerate}
An element $g \in \widetilde{PSL_2(\mathbb{R})}$ is \emph{elliptic, parabolic, or hyperbolic} if it descends to a matrix with the corresponding adjective in $SL_2(\mathbb{R})$.
\end{definition}

In order to capture the relevant information of $\widetilde{PSL_2(\mathbb{R})}$ representations, \cite{cullerdunfield} defined the \emph{translation extension locus} which captures $\widetilde{PSL_2(\mathbb{R})}$ representations which are elliptic or parabolic on the boundary, and \cite{gao} defined the \emph{holonomy extension locus}, which captures $\widetilde{PSL_2(\mathbb{R})}$ representations which are hyperbolic or parabolic on the boundary. Roughly, these capture the representation data when restricted to the boundary, making them ripe for Dehn filling considerations. In this paper, we will focus on the holonomy extension locus. 

\medskip

Suppose we are given a representation $\widetilde{\rho}: \pi_(M) \to \widetilde{PSL_2(\mathbb{R})}$ be a lift of $\rho: \pi_1(M) \to SL_2(\mathbb{R})$ such that $\rho|_{\pi_1(\partial M)}$ is either hyperbolic or parabolic. Let $\pi_1(\partial M) = \langle m, \ell \rangle$, and suppose further that $\text{trans}(\widetilde{\rho}(m)) = \text{trans}(\widetilde{\rho}(\ell)) = 0$. Then $\rho|_{\pi_1(\partial M)}$ has at least one fixed vector $v \in \mathbb{C}P^1 = \mathbb{H}^2$. Now let $\lambda_m, \lambda_\ell$ be the eigenvalues of $\rho(m)$ and $\rho(\ell)$ with respect to $v$, and take the point $(\ln(\lambda_m), \ln(\lambda_\ell)) \in \mathbb{R}^2$. The \emph{holonomy extension locus}, denote $H_{0, 0}(M)$, is the image of this map in $\mathbb{R}^2$. It is shown in \cite{gao} that $H_{0, 0}(M)$ is a locally finite union of arcs in $\mathbb{R}^2$. Crucially, Lemma 3.6 in \cite{gao} shows that real representations tending toward ideal points, when projected into the holonomy extension locus, tend toward a line through the origin with slope $-r$ in $\mathbb{R}^2$, where $r$ is the boundary slope of the surface detected by the ideal point. 

\medskip

Previous results, particularly those in \cite{gao}, identified smooth points on the character variety that have real coordinates, and then deformed these smooth points to arcs of $SL_2(\mathbb{R})$-representations. This paper will focus on deforming ideal points with real limiting character to arcs of $SL_2(\mathbb{R})$ representations, which will correspond to particular arcs on the holonomy extension locus, thus leading to intervals of orderable surgeries.

\subsection{Motivating examples}

The goal of this paper is to explain particular arcs in the holonomy extension locus which result in Dehn fillings which are half-orderable near 0. 

\begin{example}[$7_3$ knot]
The $7_3$ knot is a non-fibered genus 2 knot. Its holonomy extension locus was shown in Figure 4 in Section 4 of \cite{gao}. Indeed, $H_{0, 0}(7_3)$ contained an arc with asymptotes at slope 0 and 6, and the accompanying caption notes that no existing theorem predicts the existence of this arc. However, this arc is exactly the type of arc in the holonomy extension locus predicted by the results in this paper. SnapPy and Sage show that $7_3$ is longitudinally rigid, so any ideal point detecting the Seifert surface will lie on a component on which the trace of the longitude is non-constant. In order to fully verify that this arc can be explained by the results in this paper, it remains to show that the Seifert surface is detected by an ideal point with a real limiting character which lifts the longitude to translation number 0, and to construct an arc of $SL_2(\mathbb{R})$-representations coming from this ideal point. 
\end{example}

\begin{example}[$8_6$ knot]
The $8_6$ knot is another non-fibered genus 2 knot. In Section 8.2 of \cite{gaotwobridge}, the holonomy extension locus $H_{0, 0}(8_6)$ is shown. There is another arc in the holonomy extension locus with a horizontal asymptote with slope 0. The techniques in this paper should also predict the existence of this arc. Just as in the previous example, SnapPy and Sage show that this knot is longitudinally rigid, and this example should behave similarly to the $7_3$ knot.
\end{example}

In fact, for every known example where there is no arc in the holonomy extension locus with asymptotic slope 0, the knot is fibered. The heuristic to which this paper begins to build is the following: ``Given a nonfibered integral homology solid torus, if the limiting character at an ideal point detecting a Seifert surface is conjugate to an $SL_2(\mathbb{R})$-representation, and the trace of the homological longitude is nonconstant near that ideal point, we can reasonably expect that there is an arc in the holonomy extension locus with asymptote 0, leading to the manifold being half-orderable near 0."

\section{Main results}\label{sec:proof}

\begin{theorem}\label{thm:main}
Let $M$ be an integral homology solid torus with a nonfiber genus $g$ Seifert surface $S_g$. Suppose that $S_g$ is detected by an ideal point $x$ on a component of the character variety on which the trace of the longitude $\ell$ is non-constant, and there is an arc of representations $\rho_t: \pi_1(M) \to SL_2(\mathbb{R})$ whose characters limit toward $x$ such that for any lift $\widetilde{\rho_t}: \pi_1(M) \to \widetilde{PSL_2(\mathbb{R})}$, the translation number of $\widetilde{\rho_t}(\ell)$ is 0. Then $H_{0, 0}(M)$ admits a non-horizontal arc which has a horizontal asymptote at the $x$-axis, and $M$ is half-orderable near 0.
\end{theorem}

\begin{proof}[Proof of Theorem \ref{thm:main}]
%Note that $\partial(M \setminus S_g) \cong S_{2g}$. Consequently, \textbf{ELABORATE ON THIS} it follows that $H^2(\pi_1(M \setminus S_g)) = 0$, so any representation $\rho: \pi_1(M \setminus S_g) \to SL_2(\mathbb{R})$ lifts to $\widetilde{PSL_2(\mathbb{R})}$. Let $\widetilde{\rho_\infty} = \lim_{\text{tr}(\rho_t) \to x} \widetilde{\rho_t}|_{\pi_1(M \setminus S_g)}: \pi_1(M \setminus S_g) \to \widetilde{PSL_2(\mathbb{R})}$. Then $\widetilde{\rho_\infty}$ descends to $\rho_\infty: \pi_1(M \setminus S_g) \to SL_2(\mathbb{R})$. By assumption, $trans(\widetilde{\rho_\infty}(\ell)) = 0$. Since the translation number is continuous with respect to $\widetilde{\rho}$, we have $\text{trans}(\widetilde{\rho_t}(\ell)) = 0$ for all $t$. 
By the proof of Theorem 2.2.1 in \cite{cullershalen}, at the ideal point, the trace of the meridian approaches infinity, so near the ideal point the meridian maps to a hyperbolic matrix. Since hyperbolic matrices can only commute with hyperbolic matrices, $\rho_t$ restricted to the boundary torus is hyperbolic. Since the meridian is the generator of $H^1(M)$, we may modify the lift $\widetilde{\rho_t}$ so that the translation number of the meridian is 0. (See Section 3.4 of \cite{cullerdunfield} for more details.) We thus have an arc of $\widetilde{PSL_2(\mathbb{R})}$ representations mapping to an arc in $H_{0, 0}(M)$ that goes off to infinity, whose image in $SL_2(\mathbb{R})$ approaches an ideal point $x$ of $X(M)$ detecting the Seifert surface, i.e. a surface with boundary slope 0. By Lemma 3.6 in \cite{gao}, this arc approaches the horizontal $x$-axis as an asymptote. The arc is not the horizontal $x$-axis, since this would mean that $\text{tr}(\rho(\ell))$ was constantly equal to 2 on this arc, contradicting the assumption that $x$ is an ideal point on a component of $X(M)$ for which the trace of $\ell$ is nonconstant. Thus, there exists some $a \in \mathbb{R}$ such that for $r$ between 0 and $a$, the line through the origin with slope $-r$ passes through the arc. Since at most three Dehn fillings are reducible (Theorem 1.2 of \cite{reducible}), shrink $a$ to be small enough so that all surgeries between 0 and $a$ are irreducible. By Lemma 3.8 in \cite{gao}, this means that for all $r$ between 0 and $a$, $M(r)$ has left-orderable fundamental group, as desired. 
\end{proof}

The core difficulty in applying this theorem is to establish arcs of $SL_2(\mathbb{R})$-representations tending toward ideal points. In this paper, we will establish the existence of such arcs in the case of Seifert genus one. To do this, we first prove a key lemma.

\begin{lemma}\label{lma:offdiag}
Given 
\begin{equation}
	A = \begin{pmatrix}a & b \\ c & d \end{pmatrix} \ \ \ \ \ \Delta(A) = b - c
\end{equation}
Two parabolic or elliptic matrices $A, B \in SL_2(\mathbb{R})$ with the same trace are conjugate in $SL_2(\mathbb{R})$ if and only if $\text{sign}(\Delta(A)) = \text{sign}(\Delta(B))$. If the signs are opposite, then they are conjugate by a determinant -1 matrix in $GL_2(\mathbb{R})$.
\end{lemma}

\begin{proof}
We start with parabolic matrices. Suppose $a, b, c, d \in \mathbb{R}$ with $ad - bc \neq 0$. Notably, this means that we cannot have $a = c = 0$. For a parabolic matrix:
\begin{equation}
	M = \begin{pmatrix}a & b \\ c & d\end{pmatrix}\begin{pmatrix}1 & x \\ 0 & 1\end{pmatrix}\begin{pmatrix}d & -b \\ -c & a\end{pmatrix} = \begin{pmatrix}1 - \frac{ac}{ad-bc}x & \frac{a^2}{ad-bc}x \\ -\frac{c^2}{ad-bc}x & 1 + \frac{ac}{ad-bc}x \end{pmatrix}
\end{equation}
Then $\Delta(M) = \frac{1}{ad-bc}(a^2 + c^2)x$. The sign of $\Delta(M)$ matches the sign of $x$ if and only if $ad-bc > 0$, as desired. A similar proof works for when the trace of a parabolic matrix is -2. For an elliptic matrix:
\begin{align}
	M = \begin{pmatrix}a & b \\ c & d\end{pmatrix}&\begin{pmatrix}\cos\theta & -\sin\theta \\ \sin\theta & \cos\theta\end{pmatrix}\begin{pmatrix}d & -b \\ -c & a\end{pmatrix} \\ &= \frac{1}{ad-bc}\begin{pmatrix}(ad-bc)\cos\theta + (ac+bd)\sin\theta & -(a^2+b^2)\sin\theta \\ (c^2+d^2)\sin\theta & (ad-bc)\cos\theta - (ac+bd)\sin\theta\end{pmatrix}
\end{align}
Then $\Delta(M) = -\frac{1}{ad-bc}(a^2+b^2+c^2+d^2)\sin\theta$, whose sign is equal to $-2\sin\theta$ if and only if $ad-bc > 0$, as desired. 
\end{proof}

\begin{remark}
This proof fails for hyperbolic matrices, especially since they are conjugate to diagonal matrices, with both off-diagonal entries equal to 0. By the proof of Theorem 4.3 in \cite{goldman}, two representations $\rho_1, \rho_2: G \to SL_2(\mathbb{R})$ with the same character functions are conjugate by either a determinant +1 or -1 matrix in $SL_2(\mathbb{R})$. By the above lemma, the off-diagonal difference $\Delta$ for elliptic or parabolic matrices, \emph{but not hyperbolic matrices}, act as ``barometers" for whether or any conjugating matrix has determinant +1 or -1. This concept will be key in generating arcs of $SL_2(\mathbb{R})$-representations coming from ideal points.
\end{remark}

%\begin{lemma}
%Suppose that $\rho_1, \rho_2: F_2 \to SL_2(\mathbb{R})$ with the same trace triple, and suppose that there exists $g \in F_2$ such that $\rho_1(g), \rho_2(g)$ are conjugate elliptic matrices. Then $\rho_1$ and $\rho_2$ are conjugate in $SL_2(\mathbb{R})$. 
%\end{lemma}
%
%\begin{proof}
%We know that the two matrices are conjugate by either a determinant 1 or -1 matrix, by the work in Goldman. Since $\rho_1(g)$ and $\rho_2(g)$ are elliptic matrices that rotate in the same direction, they aren't conjugate by a determinant -1 matrix. So the entire representation must be conjugate by a matrix in $SL_2(\mathbb{R})$. 
%\end{proof}

%\begin{lemma}
%Deforming two parabolics with the same shear direction gives two elliptics rotating in the same direction. 
%\end{lemma}
%
%\begin{proof}
%Suppose we have two smooth paths $A_t, B_t: [0, \epsilon) \to SL_2(\mathbb{R})$, such that $A_0, B_0$ are parabolic matrices that are conjugate. For an elliptic matrix, the sign of the difference of the off-diagonal entries records the direction of the rotation. The same applies to a parabolic. (One can see that the difference is never 0 for an elliptic or a parabolic matrix.) This smoothly varies with the matrix, so the shear sign is preserved for both $A_t$ and $B_t$. Thus, $A_t$ and $B_t$ are elliptic matrices are rotating in the same direction. 
%\end{proof}

We now introduce terminology and notation that will be prominent in the construction of $SL_2(\mathbb{R})$-arcs tending toward ideal points. 

\medskip

Let $F_2 = \langle a, b \rangle$ be the free group on two generators. Take a representation $\rho: \pi_1(M) \to SL_2(\mathbb{C})$, and let $x = \text{tr}(\rho(a)), y = \text{tr}(\rho(b)), z = \text{tr}(\rho(ab))$. Given a word $w \in F_2$, the trace of $\rho(w)$ can be written as a polynomial in $x, y, z$ by applying the relation
\begin{equation}
	\text{tr}(X)\text{tr}(Y) = \text{tr}(XY) + \text{tr}(XY^{-1}) \ \ \ \ \ X, Y \in SL_2(\mathbb{C})
\end{equation}
For shorthand, we will denote this polynomial $\text{tr}(w)$. 

\medskip

%\begin{lemma}\label{lma:coord}
%Let $\langle m, \ell \rangle \subset F_2$ be a rank-2 free subgroup. Now pick a smooth point $P$ on a curve $C \subset X(F_2) = \mathbb{C}^3$ corresponding to the trace of a representation $\rho: F_2 \to SL_2(\mathbb{C})$ where $\rho(m), \rho(\ell), \rho(m\ell)$ all have trace $\pm 2$. Suppose that $x = \text{tr}(m), y = \text{tr}(\ell), z = \text{tr}(m\ell)$ is a set of local coordinates around $P \in \mathbb{C}^3$. Then $\text{tr}(m^p\ell^q)$ is a local coordinate at $P \in C$ for all $(p, q) \in \mathbb{Z}$ but a finite union of lines in $\mathbb{Z}^2$ and a finite set. 
%\end{lemma}
%
%\begin{proof}
%Let $U_n(t)$ be the polynomial defined by the recurrence
%\begin{equation}
%	U_0(t) = 1 \ \ \ \ \ U_{-1}(t) = 0 \ \ \ \ \ U_n(t) = tU_{n-1}(t) - U_{n-2}(t)
%\end{equation}
%So $U_n(2) = n + 1$, and $U_n(-2) = (-1)^n(n+1)$. For the derivative, we also have the recurrence
%\begin{equation}
%	U_0'(t) = U_{-1}'(t) = 0 \ \ \ \ \ U_n'(t) = U_{n-1}(t) + tU_{n-1}'(t) - U_{n-2}'(t)
%\end{equation}
%For $t = \pm 2$, we get the recurrence values 
%\begin{equation}
%	U_n'(2) = n + 2U_{n-1}'(2) - U_{n-2}'(2) \ \ \ \ \ U_n'(-2) = (-1)^{n+1}n - 2U_{n-1}'(-2) - U_{n-2}'(-2)
%\end{equation}
%Solving the recurrence, we find that
%\begin{equation}
%	U_n'(2) = \frac{n(n+1)(n+2)}{6} \ \ \ \ \ U_n'(-2) = (-1)^{n+1}\frac{n(n+1)(n+2)}{6}
%\end{equation}
%By Theorem 1 of \cite{traina}, we have 
%\begin{equation}
%	\text{tr}(m^p\ell^q)(x, y, z) = M_{1}(x, y)z + M_0(x, y)
%\end{equation}
%where $M_1(x, y), M_0(x, y)$ are sums of products of some $U_p(x), U_q(y)$. Then for all but finitely many $(p, q)$ we have $\nabla\text{tr}(m_1^p\ell_1^q)(P) = (f(p, q), g(p, q), h(p, q))$ where $f, g, h$ are nonlinear polynomials in the variables $p, q$. If the tangent hyperplane is given by $\alpha x + \beta y + \gamma z = c$, let $F(p, q) = \alpha f(p, q) + \beta g(p, q) + \gamma h(p, q)$. Then $F(p, q)$ is a nonlinear polynomial in $(p, q)$. The zeroes of $F - c$, i.e. the $(p, q) \in \mathbb{Z}^2$ for which $\text{tr}(m_1^p\ell_1^q)$ fails to be a local coordinate, are the union of a finite set of lines (coming from the linear factors of $F$) and a finite set (coming from the nonlinear factors of $F$). 
%\end{proof}

%\begin{remark}
%In most cases, this lemma holds for all but finitely many $(p, q)$, but there may be scenarios where the exceptional set includes a finite set of lines. Thankfully, this suffices for our purposes. 
%\end{remark}

The \emph{free Seifert genus} of a knot is the minimum genus of a free Seifert surface, i.e. a surface whose complement is a handlebody. In particular, the complement of a genus one free Seifert surface is a genus two handlebody $H$, whose fundamental group is isomorphic to the free group on two generators, denoted $\pi_1(H) = \langle a, b \rangle$. Then the character variety is $\mathbb{C}^3$, with coordinates given by $\text{tr}(a), \text{tr}(b), \text{tr}(ab)$. 

\begin{definition}
Given a group $G$, $\varphi: A_1 \to A_2$ an isomorphism between two subgroups, the \emph{HNN extension of $G$ with respect to $\varphi$} is $G$, with an additional group element $t$ such that for all $a \in A_1$, $tat^{-1} = \varphi(a)$. 
\end{definition} 

One hopes to construct $SL_2(\mathbb{R})$-representations by expressing $\pi_1(M)$ as an HNN extension of a gluing map between two rank-2 free subgroups of $\pi_1(H)$ corresponding to the cut punctured tori. However, if we take a representation $\pi_1(H) \to SL_2(\mathbb{R})$, it is unclear that the conjugating matrix between these two rank-2 free subgroups lies in $SL_2(\mathbb{R})$ or not. This is a subtle difference between $SL_2(\mathbb{R})$ and $SL_2(\mathbb{C})$; for two irreducible representations in $SL_2(\mathbb{C})$, the character function alone is enough to determine the conjugacy class. For irreducible $SL_2(\mathbb{R})$-representations, there are two conjugacy classes with the same character function, differing by a $GL_2(\mathbb{R})$ conjugator with determinant -1. (In $SL_2(\mathbb{C})$, an equivalent conjugation is by a matrix whose entries are all pure imaginary.) If the translation number of the longitude were nonzero, the sign of the translation number would suffice to differentiate the $SL_2(\mathbb{R})$ conjugacy classes of representations. However, we are in the case where the translation number of the longitude is zero, making the situation more subtle. The main result of this section is that in the simplest situation, we can still upgrade an $SL_2(\mathbb{R})$-representation on $\pi_1(H)$ to an $SL_2(\mathbb{R})$-representation on the HNN extension $\pi_1(M)$. 

\begin{definition}
Let $M'$ be either a hyperbolic 3-manifold with a split real place or Seifert-fibred over a hyperbolic 2-orbifold  If $M'$ is Seifert-fibered over a hyperbolic 2-orbifold $\mathcal{O}$, we say that a \emph{real holonomy representation of $M'$} is a lift to $SL_2(\mathbb{R})$ of the representation to $PSL_2(\mathbb{R})$ which is the identity on the regular fiber and equal to the holonomy representation on $\mathcal{O}$. If $M'$ is a hyperbolic 3-manifold with a real place, a \emph{real holonomy representation} is a real Galois conjugate of its holonomy representation, which goes into $SL_2(\mathbb{R})$.
\end{definition}

Recall the following setup. Suppose $M$ is an integral homology solid torus with a nonfiber free genus one Seifert surface $S$, and that $M(0)$ has only one JSJ torus. Then $M \setminus S$ is homeomorphic to a genus two handlebody. Under the Dehn filling $M(0)$, the JSJ component $M'$ is a 3-manifold with two torus boundary components. This manifold $M'$ can be obtained form $M \setminus S$ by attaching a 2-handle to the curve corresponding to the homological longitude of $M$, i.e. the boundary component of $S$. This induces a quotient map $\pi_1(M \setminus S) \twoheadrightarrow \pi_1(M')$, which kills that curve. Since $M$ is nonfibered, $M'$ is not a thickened torus. We will denote the two copies of $\pi_1(S)$ inside $\pi_1(M \setminus S)$ as $\langle m_1, \ell_1 \rangle$ and $\langle m_2, \ell_2 \rangle$. This means that the homological longitude of $M$ can be expressed as the commutator $[m_1, \ell_1] = [m_2, \ell_2]$. 

\begin{theorem}\label{thm:arc}
Let $M$ be an integral homology solid torus with a nonfiber free genus one Seifert surface $S$. Suppose that $M(0)$ has only one JSJ torus, and that its JSJ complementary region $M'$ is not a hyperbolic 3-manifold with no real places. Let $\rho: \pi_1(M \setminus S) \twoheadrightarrow \pi_1(M') \to SL_2(\mathbb{R})$ be a real holonomy representation of $M'$ composed with the natural quotient map, $\chi$ be the character of $\rho$, and $C \subset \mathbb{C}^3$ be the Zariski closure of the image of $r: X(M) \to \mathbb{C}^3$. If $\varphi$ is the gluing map between once-punctured tori on $M \setminus S$ to create $M$, say it maps $m_1$ to $m_2$ and $\ell_1$ to $\ell_2$, so $C$ is the curve defined by $\text{tr}(m_1) = \text{tr}(m_2), \text{tr}(\ell_1) = \text{tr}(\ell_2), \text{tr}(m_1\ell_1) = \text{tr}(m_2\ell_2)$. Suppose $\chi \in C$ is a smooth point, there exists some word $w$ in $m_1, \ell_1$ such that $\Delta(\rho(w)) = \Delta(\varphi(\rho(w)))$ and $\text{tr}(w)$ is a local coordinate on $C$, and that $\text{tr}([m_1, \ell_1])$ is not constantly equal to 2 near $\chi$. Then:
\begin{enumerate}
	\item $S$ is detected by an ideal point on a component of $X(M)$ for which the trace of the longitude is nonconstant, and the limiting character is $\chi$.
	\item There exists an arc of representations $\rho_t: \pi_1(M \setminus S) \to SL_2(\mathbb{R})$, where $\rho_0 = \rho$, such that the conjugating matrix between $\rho_t|_{\langle m_1, \ell_1 \rangle}$ and $\rho_t|_{\langle m_2, \ell_2 \rangle}$ lies in $SL_2(\mathbb{R})$
	\item Tor all lifts $\widetilde{\rho_t}$ to $\widetilde{PSL_2(\mathbb{R})}$, $\text{trans}(\widetilde{\rho_t}([m_1, \ell_1])) = 0$
\end{enumerate} 
\end{theorem}

\begin{proof}
The first assertion follows from Theorem \ref{thm:genusone}. To prove the second assertion, notice that $\rho(w)$ is conjugate to $\rho(\varphi(w))$ by Lemma \ref{lma:offdiag}.
%we show that there exists some parabolic element $p \in \pi_1(T_1)$ such that $\rho(p)$ and $\rho(\varphi(p))$ are conjugate parabolic matrices. There exist $T_1, T_2 \in SL_2(\mathbb{R})$ such that
%\begin{equation}
%	T_1\rho(m_1)T_1^{-1} = \begin{pmatrix}\pm 1 & a \\ 0 & \pm 1\end{pmatrix} \ \ \ T_1\rho(\ell_1)T_1^{-1} = \begin{pmatrix}\pm 1 & b \\ 0 & \pm 1\end{pmatrix}
%\end{equation}
%\begin{equation} 
%	T_2\rho(m_2)T_2^{-1} = \begin{pmatrix}\pm 1 & c \\ 0 & \pm 1\end{pmatrix} \ \ \ T_2\rho(\ell_2)T_2^{-1} = \begin{pmatrix}\pm 1 & d \\ 0 & \pm 1\end{pmatrix}
%\end{equation}
%The upper right entries of $T_1\rho(m_1^p\ell_1^q)T_1^{-1}$ and $T_2\rho(m_2^p\ell_2^q)T_2^{-1}$ are $n(p, q) = \pm pa \pm qb, m(p, q) = \pm pc \pm qd$. Let $v = (a, b), w = (c, d)$, with signs adjusted if the corresponding matrices have eigenvalues $\pm 1$. By hypothesis, $v$ is not a negative real multiple of $w$. Then there exists an infinite cone of $(p, q) \in \mathbb{Z}^2$ such that the sign of $n(p, q)$ equals the sign of $m(p, q)$ and $\rho(m_1^p\ell_1^q)$ is conjugate to $\rho(m_2^p\ell_2^q)$ by Lemma \ref{lma:offdiag}. By Lemma \ref{lma:coord}, we can pick one such $(p, q)$ such that $\text{tr}(m_1^p\ell_1^q)$ is a local coordinate; let $w_1 = m_1^p\ell_1^q$, $w_2 = m_2^p\ell_2^q$. 
Since $\text{tr}(w)$ is a local coordinate on $C$ with real coordinates in $\mathbb{C}^3$, by Lemma 2.8 in \cite{cullerdunfield}, we may smoothly deform $\chi$ to the arc of traces $\chi_t$, $t \in [0, \epsilon)$ such that $|\chi_t(w)|, |\chi_t(\varphi(w))| < 2$ and $\chi_t \in C$. Since $\rho$ is either a holonomy representation of a 2-orbifold or a real conjugate of the holonomy of a hyperbolic 3-manifold, its image must contain a hyperbolic matrix. Then in a neighborhood of $\chi$, by Theorem 4.3 of \cite{goldman}, all the $\chi_t$ must lift to $\rho_t: \pi_1(M \setminus S) \to SL_2(\mathbb{R})$. Then $\rho_t(w)$ and $\rho_t(\varphi(w))$ are elliptic matrices. By assumption, the trace of $[m_1, \ell_1] \in \pi_1(M \setminus S)$, must also not be constantly equal to 2 on this arc. Note that $\rho$ is irreducible, and so the character map $\rho \mapsto \chi_\rho$ is a local homeomorphism at $\rho$. Then $\Delta(\rho_t(w))$ and $\Delta(\rho_t(\varphi(w)))$ still have the same sign due to the continuity of the deformation. By Lemma \ref{lma:offdiag}, $\rho_t(w)$ and $\rho_t(\varphi(w))$ remain conjugate. Now let $\rho_{1, t} = \rho_t|_{\langle m_1, \ell_1 \rangle}$ and $\rho_{2, t} = \rho_t|_{\langle m_2, \ell_2 \rangle}$ for $t \in (0, \epsilon)$; since the trace of $[m_1, \ell_1] = [m_2, \ell_2]$ is not constantly equal to 2, $\rho_{1, t}$ and $\rho_{2, t}$ are irreducible for $t$ small enough. Since $\rho_{1, t}$ and $\rho_{2, t}$ have the same character function and are irreducible, they are conjugate by either a determinant 1 or -1 matrix, by the proof of Theorem 4.3 in \cite{goldman}. Since $\rho(w)$ and $\rho(\varphi(w))$ are conjugate elliptic matrices, $\rho_{1, t}$ and $\rho_{2, t}$ cannot be conjugate by a determinant -1 matrix, again by Lemma \ref{lma:offdiag}. So $\rho_1$ and $\rho_2$ must be conjugate by a matrix in $SL_2(\mathbb{R})$. Then $\rho_t$ extends to an $SL_2(\mathbb{R})$-representation $\sigma_t$ of $\pi_1(M)$, which is the HNN extension of $F_2$ with $\langle m_1, \ell_1 \rangle$ glued to $\langle m_2, \ell_2 \rangle$. We finally prove that when lifting $\rho_t$ to $\widetilde{\rho_t}: \pi_1(M) \to \widetilde{PSL_2(\mathbb{R})}$, $\text{trans}(\widetilde{\rho}([m_1, \ell_1])) = 0$. By the Milnor-Wood bound stated as Proposition 6.5 of \cite{cullerdunfield}, this translation number must have translation number either -1, 0, or 1. However, by Claim 8.5 of \cite{cullerdunfield}, since $\rho([m_1, \ell_1]) = I$, for any lift $\widetilde{\rho}$, $\text{trans}(\widetilde{\rho}([m_1, \ell_1])) = 0$. Since the meridian $\mu \in \pi_1(M)$ intersects the longitude, by the proof of Theorem 2.2.1 of \cite{cullershalen}, on $X(M)$, $\text{tr}(\sigma_t(\mu))$ approaches infinity as the ideal point detecting $S$ is approached; in particular, near enough to the ideal point, $\sigma_t(\mu)$ is hyperbolic. However, only hyperbolic matrices can commute with hyperbolic matrices, so the homological longitude $[m_1, \ell_1]$ must also be mapped to a hyperbolic matrix under $\rho_t$. Since the translation number is continuous on $\widetilde{PSL_2(\mathbb{R})}$-representations and integral on hyperbolic elements of $\widetilde{PSL_2(\mathbb{R})}$, we conclude that after choosing $\epsilon$ small enough, for all $t \in [0, \epsilon)$, $\text{trans}(\widetilde{\rho_t}([m_1, \ell_1])) = 0$, as desired. 
\end{proof}

Theorems \ref{thm:main} and \ref{thm:arc} lead to the following corollary:

\begin{corollary}\label{cor:general}
Let $S$ be a nonfiber free genus one Seifert surface in an integer homology 3-sphere $M$ such that $M(0)$ has only one JSJ torus, $M'$ be the JSJ component of $M(0)$ that is not hyperbolic with no real place, $\chi$ the trace of a real holonomy of $M'$, and $\varphi: T_1 \to T_2$ the gluing map between the two torus boundary components of $M'$. Suppose that $\chi$ is a smooth point in the Zariski-closure of the image of the restriction map $r: X(M) \to X(M \setminus S)$, there exists some word $w \in \pi_1(T_1)$ such that $\Delta(\rho(w)) = \Delta(\varphi(\rho(w)))$ and $\text{tr}(w)$ is a local coordinate, and the trace of the longitude of $M$ is not constantly equal to 2 near $\chi$. Then $H_{0, 0}(M)$ admits a non-horizontal arc which has a horizontal asymptote at the $x$-axis, and $M$ is half-orderable near 0.
\end{corollary}

%%We begin by proving a technical lemma that represents a key step in the proof of Theorem \ref{thm:main}. In essence, this proof allows us to extend $SL_2(\mathbb{R})$ representations on sutured 3-manifolds to $SL_2(\mathbb{R})$ representations on glued 3-manifolds. The main obstruction to this is that while trace data alone is enough to determine the conjugacy classes of an irreducible representation to $SL_2(\mathbb{C})$, it is not enough for conjugacy classes in $SL_2(\mathbb{R})$. 
%%
%%\begin{lemma}
%%An irreducible representation of a once-punctured surface group into $SL_2(\mathbb{R})$ with integral Euler class cannot send the puncture to an elliptic matrix.
%%\end{lemma}
%%
%%\begin{proof}
%%Elliptic elements have non-integral translation number. \textbf{FINISH THIS}
%%\end{proof}
%%
%%The following comes from \textbf{CITE GOLDMAN}.
%%
%%\begin{lemma}
%%Let $\rho_1, \rho_2: F_n \to SL_2(\mathbb{R})$ be irreducible representations with the same trace data. Then they are conjugate by either a matrix in $SL_2(\mathbb{R})$ or a matrix in $GL_2(\mathbb{R})$ with determinant -1. 
%%\end{lemma}
%%
%%\begin{proof}
%%Since they have the same trace they are $SL_2(\mathbb{C})$-conjugate. Let $F_n = \langle x_1, \dots, x_n \rangle$. By Burnside's lemma, the $\mathbb{R}$-algebra generated by $\rho_1(x_i)$ and $\rho_2(x_i)$ are both the full matrix algebra $M_2(\mathbb{R})$. So any conjugating matrix between $\rho_1$ and $\rho_2$ normalizes the full matrix algebra $M_2(\mathbb{R})$. Thus there are actually only two cases, where it's determinant plus or minus one. 
%%\end{proof}
%%
%%\begin{lemma}
%%Let $S_g$ be a once-punctured genus $g$ surface with puncture $p \in \pi_1(S_g)$. Take two irreducible representations $\rho_1, \rho_2: \pi_1(S_g) \to SL_2(\mathbb{R})$ with the same character function, and such that $\rho_1(p), \rho_2(p)$ are hyperbolic. Then $\rho_1, \rho_2$ are conjugate in $SL_2(\mathbb{R})$ if and only if $\rho_1(p)$ and $\rho_2(p)$ are conjugate in $SL_2(\mathbb{R})$.
%%\end{lemma}
%%
%%\begin{proof}
%%Suppose $\rho_1(p)$ and $\rho_2(p)$ are conjugate in $SL_2(\mathbb{R})$. Since $\rho_1$ and $\rho_2$ have the same character function, either the two representations are conjugate by a determinant -1 matrix or an element of $SL_2(\mathbb{R})$. Suppose $\rho_1$ and $\rho_2$ are conjugate by a determinant -1 matrix. Then $\rho_1(p)$ and $\rho_2(p)$ are conjugate by determinant -1. But this means that $\rho_1(p)$ and $\rho_2(p)$ are hyperbolic matrices which are conjugate by both a determinant -1 and +1 matrix. This is a contradiction, since the corresponding geodesics in $\mathbb{H}^2$ cannot be related by both an orientation-reversing and orientation-preserving isometry. 
%%\end{proof}
%%
%%\begin{lemma}\label{lma:gluing}
%%Let $M$ be a 3-manifold with boundary a surface of genus $2g$. Let $S^+, S^- \subset \partial M$ be two incompressible once-punctured genus $g$ surfaces in the boundary separated by a curve $c$. Let $\varphi: \pi_1(S^+) \to \pi_1(S^-)$ be an isomorphism induced by an orientation-reversing homeomorphism from $S^+ \to S^-$. Let $M_\varphi$ be the manifold obtained by gluing $S^+$ to $S^-$ along $\varphi$, and let $\rho: \pi_1(M) \to SL_2(\mathbb{R})$ be a representation such that $\rho|_{\pi_1(S^+)}$ is irreducible and has the same trace as $\rho|_{\pi_1(S^-)}$ when mapped under $\varphi$. Then $\rho$ extends to a representation $\rho_\varphi: \pi_1(M_\varphi) \to SL_2(\mathbb{R})$. 
%%\end{lemma}
%%
%%\begin{proof}
%%First, we note that a choice of an orientation on $M$ induces a opposite orientations on $S^+$ and $S^-$, since orientations are determined by the orientation on $c$. It follows that $\rho|_{S^+}$ and $\rho|_{S^-}$ are oppositely oriented, since the bases will send $c$ to inverse hyperbolic elements. Then since $\varphi$ is orientation-reversing, it will send a basis of $\pi_1(S^+)$ to a basis for which $\rho|_{S^-}$ sends $c$ to the same hyperbolic element, translating in the same direction. The above lemma finishes the proof. 
%%\end{proof}
%
%We now prove Theorem \ref{thm:main}, the main result of this paper. 
%
%\begin{proof}[Proof of Theorem \ref{thm:main}]
%%By assumption, there exists an ideal point $x \in X(M)$ detecting $S_g$ such that its limiting character is the trace of a real representation $\rho_\infty: \pi_1(M \setminus S_g) \to SL_2(\mathbb{R})$. By Lemma 13 in \cite{tillmann}, $x$ corresponds to a particular point with real coordinates on the sutured gluing subvariety. Since the sutured gluing subvariety is defined over $\mathbb{Q}$, by the smooth suturing assumption and Proposition 2.8 of \cite{cullerdunfield}, near $x$ there exists a real-analytic arc of irreducible representations $\rho_t': \pi_1(M \setminus S_g) \to SL_2(\mathbb{R})$, all of which lie on the sutured gluing subvariety. Thus, from Lemma \ref{lma:gluing}, there is a real-analytic arc of representations $\rho_t: \pi_1(M) \to SL_2(\mathbb{R})$ near the ideal point $x$ by gluing among $S_g$. Since $M$ is an integral homology solid torus, there is no Euler class obstruction to lifting, and we may lift $\rho_t$ to an arc $\widetilde{\rho_t}: \pi_1(M) \to \widetilde{PSL_2(\mathbb{R})}$ which approaches the ideal point. \textbf{ARGUE THAT THESE ARE HYPERBOLIC ON THE BOUNDARY.}
%%	
%%\medskip
%	
%Let $\widetilde{\rho_\infty}: \pi_1(M \setminus S_g) \to \widetilde{PSL_2(\mathbb{R})}$ be the limit of $\widetilde{\rho_t}|_{M \setminus S_g}$. By construction, $\widetilde{\rho_\infty}$ descends to $\rho_\infty: \pi_1(M \setminus S_g) \to SL_2(\mathbb{R})$. By assumption, the limiting character has Euler class 0 when restricted to $S_g$, and it follows that $trans(\widetilde{\rho_\infty}(\gamma)) = 0$. Since the translation number is continuous with respect to $\widetilde{\rho}$, we have $trans(\widetilde{\rho_t}(\gamma)) = 0$ for all $t$. 
%
%\medskip
%	
%Modify $\widetilde{\rho_t}$ so that the translation number of the meridian is 0. We thus have an arc of $\widetilde{PSL_2(\mathbb{R})}$ representations mapping into $H_{0, 0}(M)$ that goes off to infinity, approaching an ideal point detecting a surface with boundary slope 0. By the proof of Lemma 3.6 in \cite{gao}, this arc approaches the horizontal $x$-axis as an asymptote. The arc is not the horizontal $x$-axis, since this would mean that $\text{tr}(\rho(\gamma))$ was constantly equal to 2 on this arc, but $x$ is an ideal point on a component of $X(M)$ for which the trace of the longitude is nonconstant. Thus there exists some $a \in \mathbb{R}$ such that for $r$ between 0 and $a$, the line through the origin with slope $-r$ passes through the arc. By Lemma 3.8 in \cite{gao}, this means that if $M(r)$ is irreducible, it has left-orderable fundamental group, as desired. 
%\end{proof}
%
%%\begin{proof}[Proof of Corollary \ref{cor:genusone}]
%%We are in the situation of Theorem \ref{thm:genusone}, so the limiting character is the holonomy of the JSJ component of $M(0)$. Since the limiting character is always reducible, it has Euler class 0 in this case. If the JSJ component is Seifert-fibered, the representation must be to $SL_2(\mathbb{R})$. If the JSJ component is hyperbolic with real places, the image of the holonomy representation contains parabolics, meaning that the quaternion algebra generated by the image is the entire matrix algebra. This means that the representation is actually conjugate into $PSL_2(\mathbb{R})$ instead of $SU(2)$. (For more on this, see \textbf{MACLACHLAN REID}.) In either case, this holonomy representation can be conjugated to an $SL_2(\mathbb{R})$ representation. Thus we are in the situation of Theorem \ref{thm:main}, so $M$ is orderable near 0.
%%\end{proof}

\section{Examples}\label{sec:examples}

As an explicit example of Corollary \ref{cor:general} in action, we turn to odd classical pretzel knots. The $(2p+1, 2q+1, 2r+1)$ pretzel knots have free Seifert genus one, making them computationally easier to work with. Let $M_{2p+1, 2q+1, 2r+1}$ be the complement in $S^3$. Fundamental groups and JSJ decompositions of 0-surgeries on odd pretzel knots are well-understood, which is captured in the following theorem of \cite{sekino}:

\begin{theorem}[\cite{sekino}]
The group $\pi_1(M_{2p+1, 2q+1, 2r+1})$ can be written down as
\begin{equation}
	\pi_1(M_{2p+1, 2q+1, 2r+1}) = \langle a, b, t \mid ta^{r+1}(ba)^qbt^{-1} = a^{r+1}(ba)^q, tb^{p+1}(ab)^qt^{-1} = b^{p+1}(ab)^qa \rangle
\end{equation}
where $a, b$ are generators of the genus-two handlebody obtained by taking the complement of the genus 1 Seifert surface. (This is called the \emph{Lin presentation} of the fundamental group.) We have the following descriptions of JSJ decompositions of $M_{2p+1, 2q+1, 2r+1}(0)$. 
\begin{itemize}
	\item $M_{-3, 3, 2n+1}(0)$ has one JSJ component which is the $(2, 4)$ torus link, which is Seifert fibered over the annulus with a cone point of order 2.
	\item $M_{-3, 5, 5}(0)$ and $M_{3, -5, -5}(0)$ has two JSJ components: the trefoil knot complement and the trivial circle fiber over the thrice-punctured sphere, denoted $S^1 \times S_{0, 3}$. 
	\item For any other $(p, q, r)$, $M_{2p+1, 2q+1, 2r+1}(0)$ has one JSJ component which is a hyperbolic 3-manifold. 
\end{itemize}
\end{theorem}

\begin{remark}
The JSJ decomposition of $M_{-3, 5, 5}(0)$ and $M_{3, -5, -5}(0)$ suggests that there is a separating twice-punctured torus with slope 0. This becomes the torus separating the trefoil knot complement from $S^1 \times S_{0, 3}$, while the Seifert surface caps off to the essential torus bounding the two other torus boundary components of $S^1 \times S_{0, 3}$. In \cite{twicepunctured}, it is shown that the union of the Seifert-surface and the twice-punctured torus is detected by an ideal point, where the limiting character at that ideal point restricts to the holonomy representation of the thrice-punctured sphere and the $(2, 3, \infty)$ triangle orbifold underlying the trefoil knot complement. 
\end{remark}

%\begin{conjecture}
%The union of the Seifert surface and the twice-punctured torus in the $(-3, 5, 5)$ and $(3, -5, 5)$ pretzel knots is also detected by an ideal point. The limiting character restricts to the holonomy representation of the thrice-punctured sphere and the $(2, 3, \infty)$ triangle orbifold underlying the trefoil knot complement. 
%\end{conjecture}

The following lemma establishes all the necessary comptuational facts about the limiting character of an ideal point on $X(M_{-3, 3, 2n+1})$ detecting the Seifert surface. This also serves as an explicit demonstration of the proof of Theorem \ref{thm:arc} which generates arcs of $SL_2(\mathbb{R})$-representations coming from ideal points on character varieties.

\begin{lemma}\label{lma:limiting}
Let $\rho_n: F_2 \to SL_2(\mathbb{R})$ be the representation 
\begin{equation}
	\rho_n(a) = \begin{pmatrix}-1 & 1 \\ 0 & -1\end{pmatrix} \ \ \ \ \ \rho_n(b) = \begin{pmatrix}2n+1 & n \\ 2 & 1\end{pmatrix}
\end{equation}
with $\chi_n = (-2, 2n+2, -2n)$ its character. Let 
\begin{equation}
	m_1 = a^{n+1}bab \ \ \ \ \ m_2 = a^{n+1}ba \ \ \ \ \ \ell_1 = b^{-1}ab \ \ \ \ \ \ell_2 = b^{-1}aba
\end{equation}
and define a curve
\begin{equation}
	C_n = \{\text{tr}(m_1) = \text{tr}(m_2), \text{tr}(\ell_1) = \text{tr}(\ell_2), \text{tr}(m_1\ell_1) = \text{tr}(m_2\ell_2)\} \subset \mathbb{C}^3
\end{equation}
Then the following is true.
\begin{enumerate}
	\item $\rho_n(m_1)$ and $\rho_n(m_2)$ are conjugate parabolic matrices.
	\item $\chi_n \in C_n$. 
	\item $\chi_n$ is a smooth point on $C_n$.
	\item $\text{tr}(m_1)$ is a local coordinate on $C_n$.
	\item The trace of the longitude is non-constant at $\chi_n$.
	\item $\chi_n$ corresponds to an ideal point $X(M_{-3, 3, 2n+1})$ detecting the Seifert surface, and this ideal point lies on a component of $X(M)$ on which the trace of the longitude is nonconstant.
\end{enumerate}
\end{lemma}

\begin{proof}
A computation yields
\begin{equation}
	\rho_n(m_1) = (-1)^n\begin{pmatrix}-2n-1 & -n \\ 4n & 2n-1\end{pmatrix} \ \ \ \ \ \rho_n(m_2) = (-1)^n\begin{pmatrix}-1 & 0 \\ 2 & -1\end{pmatrix}
\end{equation}
Let 
\begin{equation}
	P_1 = \begin{pmatrix}1 & 0 \\ 2 & 1 \end{pmatrix} \ \ \ \ \ P_2 = \begin{pmatrix} 1 & -1 \\ 1 & 0 \end{pmatrix}
\end{equation}
\begin{equation}
	P_1\rho_n(m_1)P_1^{-1} = (-1)^n\begin{pmatrix}-1 & -n \\ 0 & -1\end{pmatrix} \ \ \ \ \ P_2\rho_n(m_2)P_2^{-1} = (-1)^n\begin{pmatrix}-1 & -2 \\ 0 & -1\end{pmatrix} 
\end{equation}
These matrices will always have the same sign of the upper-right entry, proving \textbf{(1)}. We also have 
\begin{equation}
	\rho_n(\ell_1) = \begin{pmatrix}1 & 1 \\ -4 & -3\end{pmatrix} \ \ \ \ \ \rho_n(\ell_2) = \begin{pmatrix}-1 & 0 \\ 4 & -1\end{pmatrix}
\end{equation}
Notice that $\rho_n(\ell_1)$ and $\rho_n(\ell_2)$ aren't conjugate. Finally,
\begin{equation}
	\rho_n(m_1\ell_1) = (-1)^n\begin{pmatrix}2n-1 & n-1 \\ -4n & 3-2n\end{pmatrix} \ \ \ \ \ \rho_n(m_2\ell_2) = (-1)^n\begin{pmatrix}1 & 0 \\ -6 & 1\end{pmatrix}
\end{equation}
Then $\rho_n$ satisfies the equations $\text{tr}(m_1) = \text{tr}(m_2), \text{tr}(\ell_1) = \text{tr}(\ell_2), \text{tr}(m_1\ell_1) = \text{tr}(m_2\ell_2)$, meaning that $\chi_n = (-2, 2n+2, -2n) \in C_n$. The Jacobian at $\chi_n$ is given by 
\begin{equation}
	\begin{pmatrix}(-1)^{n+1}\frac{4n^4 + 4n^3 - 7n^2 - 4n}{3} & (-1)^{n+1}(2n^2 - n - 1) & (-1)^{n+1}(2n^2 + n - 2) \\ (2n+1)^2 & 4 & 4 \\ (-1)^{n+1}\frac{-4n^4 + 4n^3 + 43n^2 + 26n + 3}{3} & (-1)^{n+1}(-2n^2 + 5n + 9) & (-1)^{n+1}(-2n^2 + 3n + 14) \end{pmatrix}
\end{equation}
The determinant of this matrix is zero, and the top left $2 \times 2$ minor is nonzero, so the rank is 2. This confirms that $\chi_n$ is a smooth point on $C_n$, proving \textbf{(3)}. In order to show that $\text{tr}(m_2)$, we show that the gradient of $\text{tr}(m_2)$ is not in the span of the gradient. We compute that the gradient is 
\begin{equation}
	\nabla \text{tr}(m_2)|_{\chi_n} = (-1)^{n+1}\left(\frac{2n(n+1)(n+2)}{3}, n+1, n+2\right)
\end{equation}
which can be seen not to be in the span of the Jacobian of $C_n$ at $\chi_n$. Thus, $\text{tr}(m_2)$ is a local coordinate at $\chi_n$, proving \textbf{(4)}. An integral basis vector for the kernel of the Jacobian is $(12, -(4n^3 + 12n^2 + 17n + 6), 4n^3 + 5n + 3)$. The Hessian of $\text{tr}([m_1, \ell_1])$ at $(-2, 2n+2, -2n)$ is given by 
\begin{equation}
	\begin{pmatrix}\frac{8n^2(n+1)^2(2n+1)^2}{9} & \frac{8n^2(n+1)(2n+1)}{3} & \frac{8n(n+1)^2(2n+1)}{3} \\ \frac{8n^2(n+1)(2n+1)}{3} & 8n^2 & 8n(n+1) \\ \frac{8n(n+1)^2(2n+1)}{3} & 8n(n+1) & 8(n+1)^2\end{pmatrix}
\end{equation}
Since this is nonzero on the tangent plane, it follows that the trace of $[m_1, \ell_1]$ is nonconstant at $\chi_n$, proving \textbf{(5)}. Now, since we already established that $\rho_n(\ell_1)$ and $\rho_n(\ell_2)$ are nonconjugate, it follows that $\rho_n|_{\langle m_1, \ell_1 \rangle}$ and $\rho_n|_{\langle m_2, \ell_2 \rangle}$ are nonconjugate representations in $SL_2(\mathbb{R})$, despite having the same character. Since $\text{tr}([m_1, \ell_1])$ is not constantly equal to 2 near $\chi_n$ in a small neighborhood of $\chi_n$ on $C_n$, $\rho_n|_{\langle m_i, \ell_i \rangle}$ for $i = 1, 2$ are irreducible with the same character. So these two restricted representations are conjugate in $SL_2(\mathbb{C})$, and they glue together to form the trace of a legitimate representation on the HNN extension gluing together $\langle m_1, \ell_1 \rangle$ to $\langle m_2, \ell_2 \rangle$, which is the fundamental group of $M_{-3, 3, 2n+1}$. Since $\chi_n$ cannot glue to form a representation, the nearby characters do not converge to a legitimate representation of $\pi_1(M_{-3, 3, 2n+1})$. Thus, they must converge to an ideal point, so $\chi_n$ corresponds to an ideal point on the character variety of $M_{-3, 3, 2n+1}$, and since nearby deformations have nonconstant longitude trace, this proves \textbf{(6)}. 
\end{proof}

\begin{remark}
From the results of \cite{sekino}, we know that the JSJ component of 0-surgery on a $(-3, 3, 2n+1)$ pretzel knot is the $(2, 4)$ torus link. In this case, $\rho_n$ is a holonomy representation of the base orbifold $A^2(2)$ of the $(2, 4)$ torus link. 
\end{remark}

Then using Corollary \ref{cor:general}, we have the following explicit result.

\begin{corollary}
The holonomy extension loci of the $(-3, 3, 2n+1)$ pretzel knot complements admit a non-horizontal arc which has a horizontal asymptote at the $x$-axis. Consequently, the $(-3, 3, 2n+1)$ pretzel knot complements are half-orderable near 0.
\end{corollary}

The orderability near 0 of the $(-3, 3, 2n+1)$ pretzel knot complements does not follow directly follow from any previous results; the results of \cite{cullerdunfield} require the knot to be fibered or for the trace field of the knot complement to have real places. By the results of \cite{sekino}, none of these pretzel knots are fibered, and the first few $(-3, 3, 2n+1)$ knot complements can be shown using SnapPy and Sage to admit no real places in their trace field. The results of \cite{gao} require longitudinal rigidity, which according to SnapPy and Sage, is not satisfied by any $(-3, 3, 2n+1)$ pretzel knot for $n$ small. (However, the Alexander polynomial condition for Gao's Theorem 5.1 does seem to be satisfied by these pretzel knots, making longitudinal rigidity the only obstruction.) The results of \cite{khantran} demonstrate orderability near 0 for $(2p+1, 2q+1, 2r+1)$ pretzel knots with $p, q, r > 0$. Thus, the $(-3, 3, 2n+1)$ provide new explicit additions to the list of knot surgeries which are proven to be left-orderable. 

%\begin{remark}
%Some computations using SnapPy and Sage show that odd pretzel knots seem to never be longitudinally rigid, with at least one curve component coming from the character variety of the 0-surgery. We thus make the following conjecture.
%\end{remark}
%
%\begin{conjecture}
%Odd pretzel knots are not longitudinally rigid but are longitudinally semi-rigid.
%\end{conjecture}

%\subsection{Other examples}
%
%We first demonstrate that the manifold $M$ in Theorem \ref{thm:main} need not be an integral homology solid torus, as long as the Euler class of the arc of representations is still trivial.
%
%\begin{example}[$m006$]
%	This manifold has a genus one Seifert surface which was shown to be detected by an ideal point in \cite{dfj}. \textbf{COMPUTE EULER CLASS}. SnapPy and Sage compute that $m006$ is longitudinally rigid. It remains to show cutting along the Seifert surface results in a smoothly sutured handlebody. 
%\end{example}
%
%\begin{example}[$m037$]
%	This manifold has a genus one Seifert surface which was shown to be detected by an ideal point in \cite{dfj}. \textbf{COMPUTE EULER CLASS}. SnapPy and Sage compute that $m037$ is longitudinally rigid. It remains to show cutting along the Seifert surface results in a smoothly sutured handlebody. 
%\end{example}

%\subsection{$SU(2)$ representations}
%
%One major roadblock in establishing arcs of $SL_2(\mathbb{R})$ representations comes from orientation issues. In short, two representations having the same trace triple does not necessarily imply conjugacy. This difficulty is exacerbated by the representation being Euler class 0. This is a major contrast to the $SL_2(\mathbb{C})$ case, and as we will see, $SU(2)$.
%
%\begin{lemma}
%Let $\rho_1, \rho_2: G \to SU(2)$. If $\chi_{\rho_1}(g) = \chi_{\rho_2}(g)$ for all $g \in G$, then $\rho_1$ and $\rho_2$ are conjugate. 
%\end{lemma}
%
%\begin{proof}
%
%\end{proof}
%
%Since arcs of $SU(2)$ representations are of general interest, we apply the techniques in this paper to more easily generate arcs of $SU(2)$ representations

\section{Future directions}\label{sec:future}

In this section we discuss various methods to improve the applicability and scope of Theorem \ref{thm:main}. 

%The main obstruction to this is establishing arcs of boundary-hyperbolic representations coming from ideal points detecting Seifert surfaces. Our results also prove that many 3-manifolds with detected Seifert surface have an arc in the holonomy extension locus with a horizontal asymptote at the $x$-axis, so we are interested in determining which slope this arc in the holonomy extension locus ends. 

%\subsection{Verifying the hypotheses of Theorem \ref{thm:main}}
%
%The phenomenon where curves corresponding to character varieties of 0-surgeries collapse to a point under the restriction map is worth additional exploration. We have the following illuminating example.
%
%\begin{prop}
%The curve of abelian representations collapses to a point under the restriction map. 
%\end{prop}
%
%\begin{proof}
%The abelian representations factor through $H_1(M)$ \textbf{FINISH THIS}
%\end{proof}
%
%This could be the mechanism through which components containing irreducible representations also disappear under the restriction map, which would explain why odd pretzel knots appear to be longitudinally semi-rigid but not longitudinally rigid. Indeed, it specifically appears as though odd pretzel knot complements $M$ contain curve components coming from $M(0)$, which collapse to a point under the restriction map. We thus conjecture the following.
%
%\begin{conjecture}
%Any curve on $M$ that comes from a curve on $M(0)$ is collapsed to a point under the restriction map. 
%\end{conjecture}
%
%This, combined with results of Boyer-Zhang on non-norm curves, would prove that all knot complements are longitudinally semi-rigid.
%
%\medskip
%
%For a more theoretical approach to prove smooth suturing, one would want to establish local coordinates for the gluing subvariety. This is certainly the most approachable in the case of a free genus 1 Seifert surface, since that has $\mathbb{C}^3$ as its character variety. We have the following conjecture.
%
%\begin{conjecture}
%Let $x$ be a point corresponding to a representation killing the central annulus in a genus 2 sutured handlebody. Then the local coordinates around $x$ are given by \textbf{FIGURE THIS OUT}
%\end{conjecture}
%
%If this conjecture is proven, then one could utilize these local coordinates to show smooth suturing. This approach is taken in \cite{paoluzziporti}. Assuming this conjecture, one would also need to write down the gluing equations in terms of the local coordinates to establish smoothness. 

\subsection{Other odd pretzel knots}

We know, from the work of \cite{sekino}, most other odd pretzel knots yield 2-cusped hyperbolic 3-manifolds when taking the JSJ component of the 0-surgery. In order to establish arcs of real points coming from ideal points detecting their Seifert surfaces, we must begin with a real conjugate of the hyperbolic holonomy representation of this hyperbolic JSJ component. The author used SnapPy, Sage, and Regina to make the following computation for some other pretzel knots.

\begin{corollary}
Let $M_{p, q, r}$ be one of the following pretzel knot complements:
\begin{itemize}
	\item $(3, 3, 5), (3, 5, 5), (3, 3, 7), (3, 5, 7)$
	\item $(-5, 5, 5), (-3, 5, 9), (-3, 5, 13)$
\end{itemize} 
Then $M_{p, q, r}$ the resulting hyperbolic 3-manifold in the JSJ decomposition has a real place, and the Zariski-closure of the restriction map $r: X(M_{p, q, r}) \to \mathbb{C}^3$ is smooth. Then $H_{0, 0}(M_{p, q, r})$ admits a non-horizontal arc which has a horizontal asymptote at the $x$-axis, and $M_{p, q, r}$ is half-orderable near 0.
\end{corollary}

We make the following conjecture. 

\begin{conjecture}
Let $M$ be any odd pretzel knot complement. If the JSJ component of $M(0)$ is not a hyperbolic 3-manifold whose trace field has no real places, then $H_{0, 0}(M)$ admits a non-horizontal arc which has a horizontal asymptote at the $x$-axis, and $M$ is half-orderable near 0.
\end{conjecture}

The author observed the following interesting patterns in his computational experiments on character varieties of odd pretzel knots.
\begin{itemize}
	\item The genus of the image of canonical component appears to be equal to $n$ for the $(-3, 3, 2n+1)$ pretzel knot. This, combined with the Riemann-Hurwitz theorem, provides a lower bound for the genus of the canonical component of such pretzel knots. This could lead to effective bounds on the genera of canonical components, which are largely unknown outside of two-bridge knots \cite{petersen} and once-punctured torus bundles of tunnel number one \cite{bakerpetersen}. 
	\item The degree of the trace field of JSJ component on 0-surgery on the $(-3, 5, 2n+1)$ knot appears to be equal to $n - 1$. If this is true in general, the JSJ components of the 0-surgeries on the $(-3, 5, 4n+1)$ knot complements would have odd degree trace field, meaning they have real places, and they would be another infinite family of odd pretzel knots which are likely to be half-orderable near 0.
\end{itemize}

\subsection{Finding arcs}

The extra hypotheses in Corollary \ref{cor:general} seem reasonably generic. We conjecture the following.  

\begin{conjecture}
Let $S$ be a surface in $M$ detected by an ideal point $x$ of the character variety, and let $C \subset X(M \setminus S)$ be the Zariski-closure of the restriction map $r: X(M) \to X(M \setminus S)$. Then the limiting character at $x$ is always a smooth point in $C$ where the trace of all but finitely many words in $\pi_1(M \setminus S)$ is a local coordinate.
\end{conjecture}

These conditions all seem reasonably generic, and if proven, it would lead to a proof of the following general conjecture, with most of the extra technical conditions removed:

\begin{conjecture}
Suppose $M$ has a connected free genus one Seifert surface $S$ which is detected by an ideal point on a component of the character variety on which the trace of the longitude is nonconstant, and that $M(0)$ has only one JSJ component which is not a hyperbolic 3-manifold whose trace field has no real places. Then $H_{0, 0}(M)$ admits a non-horizontal arc which has a horizontal asymptote at the $x$-axis, and $M$ is half-orderable near 0.
\end{conjecture}

In order to further generalize these results for genus greater than one, we need to understand limiting characters of these detected Seifert surfaces. To begin, we pose the following question.

\begin{question}
What is the limiting character at an ideal point detecting a Seifert surface of genus greater than one? In particular, when lifting to $\widetilde{PSL_2(\mathbb{R})}$, is the translation number at the longitude zero?
\end{question}

Of particular interest are the $7_3$ and $8_6$ knots, whose holonomy extension loci are already known to contain the asymptote-0 arcs which motivated this paper. 

\subsection{Expanding the orderable interval}

Theorem \ref{thm:main} establishes that certain knot complements are half-orderable near 0. One would hope to obtain more precise information about the arc in the holonomy extension locus; in particular:

\begin{question}
How far does the interval of orderable slopes go? Does this interval go positive or negative?
\end{question} 

We have the following conjecture, inspired by Remark 16.23 of \cite{dunfieldrasmussen}.

\begin{conjecture}
For any Montesinos knot complement $M$, any boundary-parabolic representation into $SL_2(\mathbb{R})$ cannot have longitudinal height zero. In other words, if $\ell$ is the homological longitude of $M$, lifting a boundary-parabolic $\rho: \pi_1(M) \to SL_2(\mathbb{R})$ to $\widetilde{\rho}: \pi_1(M) \to \widetilde{PSL_2(\mathbb{R})}$ results in $\text{trans}(\widetilde{rho}(\ell)) \neq 0$. 
\end{conjecture}

%\begin{proof}
%A boundary parabolic representation restricted to $M \setminus S$ must restrict to two representations which are conjugate via a parabolic matrix, since this element is the meridian in the knot complement. This means that the representation restricted to the Seifert surface had to be parabolic to begin with. 
%\end{proof}

This conjecture serves to rule out the possibility that an arc would ``escape" the holonomy extension locus and enter the translation extension locus. We also have the following lemma, which comes from the results of \cite{Dunfield1999CyclicSD}:

\begin{lemma}
On the canonical component, arcs of $SL_2(\mathbb{R})$ representations can only end at ideal points or elliptic representations.
\end{lemma}

%\begin{proof}
%The restriction map from the character variety of the 3-manifold to the character variety restricted to the boundary is one-to-one, as proven by \textbf{DUNFIELD}. There is thus no possibility of a smooth arc leaving. The only possible ending points are singular, i.e. ideal points or elliptic representations.  \textbf{prove}
%\end{proof}

The above lemma and conjecture would show that in the situation of Corollary \ref{cor:general}, if the ideal point detecting the Seifert surface lies on the canonical component, $M$ is orderable on the interval $[0, b_1]$ or $[b_2, 0]$, where $b_1$ is the minimal positive boundary slope and $b_2$ is a maximal negative boundary slope. For example, the experimental evidence gathered for this paper suggests that the genus one Seifert surfaces of the $(-3, 3, 2n+1)$ pretzel knots are detected by an ideal point on the canonical component. It remains to see if such knots admit boundary-parabolic representations which lift the longitude to an element in $\widetilde{PSL_2(\mathbb{R})}$ with translation number 0.

\subsection{Other slopes}

The techniques in this paper become more subtle when considering incompressible surfaces whose boundary slopes are not 0. For example, for the figure-eight knot, the first example in Section 4 of \cite{gao} shows that there is an arc on the holonomy extension locus resulting in an interval of orderable slopes between -4 and 4, indicating that there is an arc of $SL_2(\mathbb{R})$-representations coming from both ideal points on its canonical component. (These ideal points detect twice-punctured tori with slope $\pm 4$.)

\medskip

However, consider the $(-2, 3, 7)$ pretzel knot. In \cite{gordonluecke}, it was shown that there is an incompressible twice-punctured torus with slope 37/2, and that 37/2-Dehn filling admits one JSJ torus. The JSJ complementary regions are fibered over the disk with two cone points. Thus, there may be an arc of $SL_2(\mathbb{R})$-representations coming from this ideal point, particularly since the issue with HNN-extensions is no longer present. In \cite{niepretzel}, it was shown that 37/2 surgery on the $(-2, 3, 7)$ pretzel knot is not left-orderable. In this case, the only obstruction to left-orderability is the translation number of the lift of the longitude, implying that the real $SL_2(\mathbb{R})$-arc at that ideal point would not lift to a $\widetilde{PSL_2(\mathbb{R})}$-representation whose longitudinal translation number is zero. These two examples lead to the following question.

\begin{question}
Let $M$ be an integral homology solid torus with homological longitude $\ell$. For which boundary slopes detected by ideal points can we establish lifts of $SL_2(\mathbb{R})$ representations lifting to $\widetilde{PSL_2(\mathbb{R})}$-representations $\widetilde{\rho_t}$ such that $\text{trans}(\widetilde{\rho_t}(\ell)) = 0$?
\end{question}

In future work, the author will explore this question for boundary slope $\infty$, particularly the case of essential Conway spheres.

\bibliography{sl2r.bib}
\bibliographystyle{plain}
	
\end{document}
